\section{The rational elements in the row-counting kernel}

% BEGIN C83N1
\paragraph{C83N1. Trace rigidity on the lattice of the actual test point.}
Use a row datum $B$ in the counting theorem, of rank $n\ge1$
over a CM extension $E/F$ with $[F:\mathbb Q]=f$. At $x$ in
the adelic unitary group, write $\Lambda_x$ for the rational
lattice cut out by its finite translate of the reference lattice.
Let $g_0>4n$ with $(g_0)$ radical in $\mathcal O_E$; this holds
at the completely split counting prime. For a rational group
element represented by $G$, suppose the kernel
$\Phi(x^{-1}Gx)$ is nonzero. Its finite idempotent is therefore
nonzero at the finite component. The two support conditions imply
\[
  G\Lambda_x=\Lambda_x,\qquad
   (G-1)^n\Lambda_x\subset g_0\Lambda_x .
\]
Indeed the conjugate fixes the reference lattice at each finite
place and its reduction is unipotent at the places over $g_0$;
conjugating back and intersecting with the rational vector space
gives both inclusions. This argument retains each original
cylinder idempotent of ARO26 on its own finite lattice.

For every integer $j\ge1$, $t_j=\operatorname{tr}(G^j)-n$
belongs to $g_0\mathcal O_E$. A full lattice over this Dedekind
ring is locally free. In local bases at all finite primes, $G$
has integral entries, so $t_j$ is integral. Modulo each prime
over $g_0$, $G^j$ is unipotent, and a nilpotent matrix has trace
zero over the residue field in every characteristic. Hence
$t_j$ lies in all these primes and in their product $(g_0)$.

Write $A_v$ for the matrix of $x^{-1}Gx$ in the archimedean
frame at $v$, and $d_v$ for the Euclidean norm of the logarithms
of its singular values. If every $d_v\le\epsilon$ and
$n\epsilon\le1/4$, then $\|A_v\|_{\rm op}\le e^\epsilon$.
For $1\le j\le n$,
$|\operatorname{tr}(A_v^j)-n|\le n(e^{n\epsilon}+1)<3n$.
These traces are embeddings of $\operatorname{tr}G^j$ by
conjugacy; their complex conjugates supply the other embeddings
of $E$. Thus every conjugate of the integer $t_j/g_0$ has
absolute value less than one. The integer norm forces $t_j=0$.
Newton's identities in characteristic zero now give the same
characteristic polynomial as the identity matrix from the first
$n$ power traces. Cayley--Hamilton gives $(G-1)^n=0$ in the
rational matrix algebra, and hence in every local realization.
% END C83N1

% BEGIN C83N2
\paragraph{C83N2. The near radius and the definite counting formula.}
Set $\epsilon=c/(n^8 f^2)$ for $0<c\le10^{-5}$ and
$\delta=3\sqrt{n\epsilon}$. Then $n\epsilon\le1/4$,
$\delta<1/4$, and $\delta qnf\le1/100$ for $q^2\le n$,
since $(\delta qnf)^2\le9c/n^4$. For each near element of
C83N1 the equality $\operatorname{tr}A_v=n$ gives
\[
 \|A_v-I\|_{\rm HS}^2
 =\operatorname{tr}(A_v^*A_v)-n
 \le n(e^{2\epsilon}-1)\le 3n\epsilon .
\]
In particular the Hilbert--Schmidt distance is at most $\delta$
in every frame. Together with the finite support and nilpotence
already proved, this gives the full original $\mathrm{NearUnip}$
at every supported rational point, including $c=10^{-6}$ as used
in $\mathrm{IsNear}$. No estimate for the size of a matrix
coefficient is needed here.

For a definite row datum every archimedean conjugate is unitary,
and has $d_v=0$. Every rational element contributing on the
kernel diagonal is therefore unipotent and unitary, hence is
the identity, since a unitary matrix is diagonalizable.
This proves the original count for the definite subcase:
\[
  D_{\rm count}=(\prod_v d(\sigma_v))\,
       \mu(\mathscr F)\operatorname{Re} e(1)=D^\star .
\]
To justify the trace identity here, convolve by the finite
self-adjoint idempotent $e$ and at each compact infinite group
by the conjugate coefficient of its lowest unit vector times
the formal degree $d(\sigma_v)$. Schur orthogonality makes
the latter the rank-one projector in the $\sigma_v$ factor
and zero on inequivalent factors. The finite and infinite
convolutions commute, so their product projects onto the
finite-level multiplicity space evaluated at the product vector.
The compactness at infinity and finiteness of definite genus
classes give finitely many double classes at any fixed compact
open subgroup. Passing to a common subgroup fixing $e$ on both
sides shows that this projector has finite rank by the compact
group decomposition. Its kernel is the unfolded sum of the
product test over rational translates. On the compact quotient
the sum is locally finite, since the lifted kernel has compact
support and the diagonal rational group is discrete. Integrating
its diagonal computes the trace, by expressing the finite-rank
orthogonal projection in an orthonormal basis. The diagonal
value is $(\prod_v d(\sigma_v))e(1)$ by the preceding argument;
$e(1)$ is real by self-adjointness. The identity uses the given
Haar measures in the coefficient orthogonality and the quotient
volume. The indefinite counting problem at the ranks occurring
in the family also has nontrivial unipotent terms.
% END C83N2

% BEGIN C83N3
\paragraph{C83N3. Separation on the actual lattice and the size of an arithmetic ball.}
Write $\tau=\min(1/2,\operatorname{least}(\Lambda_x))$,
where the least size is the infimum over nonzero sublattices
with the metrics in the frames of $x$. For a lattice $M$ of
rank $r>0$ use the normalized size
\[
   s(M)=\left(2^{fr}\operatorname{covol}(M)/
                |D_E|^{r/2}\right)^{1/(2fr)} .
\]
The largest size of a nonzero torsion-free quotient of
$\Lambda_x$ is the reciprocal of its least size. Indeed
dualizing identifies such quotient lattices with saturated
sublattices of the dual, with the subspace and quotient
metrics dual. The normalized sizes are reciprocal: in a
pseudo-basis the fractional ideals are inverted in the
dual, and the Gram determinants of dual bases are reciprocal.
Passing to saturation only decreases the covolume of a
sublattice of the same rank. The point lattice is self-dual
for its Hermitian pairing and this identification, conjugate
linear on the coefficients, is isometric in the given frames.
Positivity of the least size follows also from discreteness of
the exterior-power lattices: for each rank the determinant
ideal of a sublattice bounds a nonzero exterior vector after
one of finitely many ideal-class normalizations; the product
of its archimedean norms is bounded below by a fixed positive
constant, by selecting a nonzero coordinate and its ideal norm.

For two distinct rational elements $G,G'$ whose conjugates
at $x$ have finite component in the reference stabilizer,
the difference $A=G-G'$ induces a lattice isomorphism from
$Q=\Lambda_x/(\Lambda_x\cap\ker A)$ to its image $M\subset\Lambda_x$.
Let $a_v$ be the Hilbert--Schmidt norm of the difference in
the local frames. With the quotient norm on $Q$, each of the
$r$ singular values of the induced map at $v$ is at most $a_v$.
The Euclidean covolume changes by the product of the squared
absolute determinants at the $f$ complex places. Consequently
\[
       \tau^2\le s(M)/s(Q)\le(\prod_v a_v)^{1/f}.
\]
At a point with least size at least $1/2$ a nonidentity near
element would have product norm at most $\delta<1/4$ by C83N2,
contradicting this inequality with $G'=1$. The native
nonidentity near sum vanishes identically on this part of the
original fundamental domain, for every supplied cylinder.

More generally, if $d_v\le R_v$ with $R_v\ge0$, the matrices
have Hilbert--Schmidt norm at most $\sqrt n e^{R_v}$. Scale
the component at $v$ by $e^{\bar R-R_v}$ with
$\bar R=f^{-1}\sum_v R_v$. The product of the scale factors
is one, so the rescaled matrix tuples are separated by
$\tau^2$ in maximum norm. They lie in a ball of radius
$\sqrt n e^{\bar R}$ in real dimension $2n^2f$. Disjoint
open balls of radius $\tau^2/2$ about any finite subset give
the cardinality bound $(1+2\sqrt n e^{\bar R}/\tau^2)^{2n^2f}$.
It is independent of the chosen subset, so the ball is finite.
Since $\sqrt n e^{\bar R}/\tau^2\ge1$, this proves the native
$\mathrm{AutBallCount}$ with upper bound
\[
          A_x e^{2n^2\sum_v R_v},\qquad
          A_x=(3\sqrt n/\tau^2)^{2n^2 f}.
\]
% END C83N3

% BEGIN C83N4
\paragraph{C83N4. Absolute summation outside the small Cartan ball.}
On the same indefinite row, retain the following real
coefficient estimate with the given gap $h$ and
$1\le q\le p$ for the signature $(p,q)$:
\[
 |\Phi(x^{-1}Gx)|\le a\prod_v
 \exp\!\left[-\frac h8\min\{d_v^2,d_v/\sqrt{2q}\}\right],
 \qquad a=(\prod_v d(\sigma_v))\operatorname{Re}e(1)\ge0 .
\]
This is the pointwise decay premise on the actual lowest types.
Let $0<\epsilon\le1/\sqrt{2q}$, $2n^2\epsilon\le\log2$,
$h\epsilon\ge64n^2$ and $T=h\epsilon^2/16\ge1$.
The absolute sum over all rational elements outside
$\mathrm{IsNear}$ is at most
\[
          a A_x\,2f\,3^f e^{-T}.
\]
If $a=0$ every term vanishes. Otherwise omit zero terms and
partition each $d_v$ into $[0,\epsilon]$ and
$(2^j\epsilon,2^{j+1}\epsilon]$, $j\ge0$. Every tuple
counted here has at least one positive shell. By C83N3 the
number in a tuple is bounded by $A_x$ times the product of
the counting costs $\exp(2n^2 R_v)$ at its upper radii.
The first shell costs at most two. On a positive shell
$\min(d_v^2,d_v/\sqrt{2q})\ge\epsilon d_v$.
Its decay is at most $e^{-2T2^j}$ and its counting cost at
most $e^{T2^j}$ by $h\epsilon\ge64n^2$. The positive-shell
weights sum to at most $e^{-T}/(1-e^{-T})\le2e^{-T}$.
All shell weights at a place therefore sum to at most three.
Specify one of the $f$ positive places, sum there first, and
use the unrestricted product at the others. This proves
absolute summability and the required bound, also when several
places are positive. It is the original far-term estimate
before integrating over the fundamental domain.

In the counting range set $A=C_g\ge0$, $N=n^2f$,
$X=nf\,\mathrm{rd}\,g b$ and use $\epsilon=10^{-6}/(n^8f^2)$.
The assumptions include $q^2\le n$, $1\le g\le b^A$,
$b\ge2$, $\log f\le\sqrt n$, $\mathrm{rd}\ge1$, and
$h\ge8\cdot10^{14}(A+4)n^{19}f^5\log(b\,\mathrm{rd})$.
Thus the shell hypotheses hold for large $n$, uniformly for
these data, and $T\ge50(A+4)n^3 f\log(b\,\mathrm{rd})$.
On the moderate set $\operatorname{least}(\Lambda_x)\ge X^{-11q}$,
we have $\tau\ge X^{-11q}$ because $X^{-11q}<1/2$.
Consequently, after division by $n^3f$,
\[
 \frac{\log A_x+\log(2f)+f\log3+N}{n^3 f}
 \le 44+o(1)+\frac{44(A+1)}{\sqrt n}\log(b\,\mathrm{rd}).
\]
Indeed $\log A_x\le2N(\log(3\sqrt n)+22q\log X)$ and
$\log X\le\log n+\sqrt n+(A+1)\log(b\,\mathrm{rd})$.
The bound for $T$ dominates this quantity after the same
division, since $\log(b\,\mathrm{rd})\ge\log2$. Integration
on the moderate subset of the finite-volume fundamental domain
therefore gives at most
$e^{-n^2f}(\prod_v d(\sigma_v))\mu(\mathscr F)\operatorname{Re}e(1)$,
the exact $\mathrm{FarModerate}$ allowance. The only analytic
input in this deduction is the displayed pointwise decay,
with the source's measure and its row vector. The near terms
over the exceptional flags and the trace on the deep region
require their own estimates for an indefinite row.
% END C83N4
